\hypertarget{chapter-09-page-207}{}
\subsection{Cotlar 引理}
\label{sec:ch9-cotlar-lemma}

本章回到第~\ref{chap:singular-integrals-two}~章首先考虑的一个问题: 给定带相伴核 $K$ 的 Calderón--Zygmund 算子 $T$, $T$ 何时在 $L^2$ 上有界? 若 $T$ 是卷积算子, 问题归结为证明 $\widehat K\in L^\infty$. 对非卷积算子, 问题困难得多. 但在某些情形中, M. Cotlar 的下述结果很有用.

\begin{Lemma}[Cotlar's Lemma]
\label{lem:ch9-cotlar}
设 $H$ 是 Hilbert 空间, $\{T_j\}$ 是 $H$ 上的有界线性算子序列, $\{T_j^*\}$ 为其伴随算子, $\{a(j)\}$ 为非负数序列, 并且
\[
 \|T_i^*T_j\|_{H,H}+\|T_iT_j^*\|_{H,H}\le a(i-j).
\]
则对所有整数 $n\le m$,
\[
 \left\|\sum_{j=n}^mT_j\right\|_{H,H}
 \le\sum_{i=-\infty}^{\infty}a(i)^{1/2}.
\]
\end{Lemma}

\begin{Proof}
令
\[
 S=\sum_{j=n}^mT_j.
\]

\hypertarget{chapter-09-page-208}{}
则 $\|S\|=\|SS^*\|^{1/2}$. 事实上, 对任意整数 $k>0$,
\[
 \|S\|=\|(SS^*)^k\|^{1/(2k)}.
\]
又
\[
 (SS^*)^k=
 \sum_{j_1,\ldots,j_{2k}=n}^m
 T_{j_1}T_{j_2}^*\cdots T_{j_{2k-1}}T_{j_{2k}}^*.
\]
每个求和项的范数可用两种方式估计:
\[
 \begin{aligned}
 \|T_{j_1}T_{j_2}^*\cdots T_{j_{2k-1}}T_{j_{2k}}^*\|
 &\le\|T_{j_1}T_{j_2}^*\|\cdots
 \|T_{j_{2k-1}}T_{j_{2k}}^*\|\\
 &\le a(j_1-j_2)\cdots a(j_{2k-1}-j_{2k}),
 \end{aligned}
\]
以及
\[
 \begin{aligned}
 \|T_{j_1}T_{j_2}^*\cdots T_{j_{2k-1}}T_{j_{2k}}^*\|
 &\le\|T_{j_1}\|\,\|T_{j_2}^*T_{j_3}\|\cdots
 \|T_{j_{2k-2}}^*T_{j_{2k-1}}\|\,\|T_{j_{2k}}^*\|\\
 &\le a(0)^{1/2}a(j_2-j_3)\cdots
 a(j_{2k-2}-j_{2k-1})a(0)^{1/2}.
 \end{aligned}
\]
取两个上界的几何平均并对所有 $j_i$ 求和, 得
\[
 \begin{aligned}
 \|(SS^*)^k\|
 \le a(0)^{1/2}\sum_{j_1,\ldots,j_{2k}=n}^m
 &a(j_1-j_2)^{1/2}a(j_2-j_3)^{1/2}\cdots\\
 &\cdots a(j_{2k-1}-j_{2k})^{1/2}.
 \end{aligned}
\]
固定 $j_1,\ldots,j_{2k-1}$ 并先对 $j_{2k}$ 求和, 再依次对 $j_{2k-1}$ 直至 $j_1$ 求和. 于是
\[
 \|(SS^*)^k\|
 \le a(0)^{1/2}(m-n+1)
 \left(\sum_{i=-\infty}^{\infty}a(i)^{1/2}\right)^{2k-1},
\]
其中 $m-n+1$ 是对 $j_1$ 求和时出现的项数. 因此
\[
 \|S\|=\|(SS^*)^k\|^{1/(2k)}
 \le(m-n+1)^{1/(2k)}
 \sum_{i=-\infty}^{\infty}a(i)^{1/2}.
\]
令 $k\to\infty$ 即得所需结论.
\end{Proof}

若对 $H$ 的一个稠密子集中的每个 $x$, 和式 $\sum_{j=-\infty}^{\infty}T_jx$ 存在, 则可将其延拓为 $H$ 上的算子, 且算子范数由 $\sum_i a(i)^{1/2}$ 控制. 事实上, 无需假设该级数在稠密子集上存在, 因为 Cotlar 引理的假设已经蕴含它对每个 $x\in H$ 收敛.

\hypertarget{chapter-09-page-209}{}
\begin{Example}[An application to the Hilbert transform]
\label{ex:ch9-cotlar-hilbert-transform}
令 $H=L^2(\mathbb R)$, 并对 $f\in L^2(\mathbb R)$ 定义
\[
 T_jf(x)=\int_{2^j<|t|\le2^{j+1}}f(x-t)\,\frac{dt}{t}.
\]
则 $T_j$ 的任意有限和在 $L^2$ 上一致有界, 因而 Hilbert 变换的截断积分一致有界.
\end{Example}

对每个整数 $j$, $|T_jf(x)|\le4Mf(x)$, 其中 $M$ 是 Hardy--Littlewood 极大算子, 所以 $T_j$ 显然在 $L^2(\mathbb R)$ 上有界. 用 Cotlar 引理证明任意有限和一致有界. 严格说来, 论证先得到在二进数 $2^j$ 处截断的积分有界, 但很容易推广到任意其他截断.

因为 $T_j^*=-T_j$, 只需估计 $\|T_iT_j\|$. 令
\[
 K_j(x)=x^{-1}\chi_{\Delta_j}(x),
 \qquad
 \Delta_j=\{x\in\mathbb R:2^j<|x|<2^{j+1}\}.
\]
则 $T_jf=K_j*f$, $T_iT_jf=K_i*K_j*f$, 因而
\[
 \|T_iT_j\|\le\|K_i*K_j\|_1.
\]
又
\[
 K_i*K_j(x)=\int_{\mathbb R}\frac1t\chi_{\Delta_i}(t)
 \frac1{x-t}\chi_{\Delta_j}(x-t)\,dt.
\]
不妨设 $i<j$ 且 $x>0$. 若
\[
 x\notin(2^j-2^{i+1},2^{j+1}+2^{i+1}),
\]
则 $K_i*K_j(x)=0$. 在该区间上 $|K_i*K_j(x)|\le2\cdot2^{-j}$. 对区间
\[
 (2^j-2^{i+1},2^j+2^{i+1}),
 \qquad
 (2^{j+1}-2^{i+1},2^{j+1}+2^{i+1})
\]
使用这个估计. 在其余部分, 若 $t\in\Delta_i$, 则 $x-t\in\Delta_j$, 故
\[
 K_i*K_j(x)=\int_{\mathbb R}\frac1t\chi_{\Delta_i}(t)
 \left(\frac1{x-t}-\frac1x\right)\chi_{\Delta_j}(x-t)\,dt.
\]
因此 $|K_i*K_j(x)|\le2\cdot2^i2^{-2j}$, 从而
\[
 \|K_i*K_j\|_1\le C2^{-|i-j|}.
\]
这个估计使我们可以按所需方式应用 Cotlar 引理.
